\begin{frame}{확인 문제: 장을 넘나드는 복습 4개 (1)-(2)}
  \small
  \begin{itemize}
    \item \textbf{1.} Ch2.2의 \kb{UCB}와 Ch15.2의 \kb{UCT}는 수식이
      거의 똑같다. 무엇이 같고, 다른가?
      \vskip0.2em
      {\footnotesize $\to$ 둘 다 ``지금까지의 평균 성과 + 덜
      시도해본 정도에 대한 불확실성 보너스''라는 같은 형태
      ($\bar{X} + c\sqrt{\ln N / n}$). \textbf{다른 것은 적용 대상} —
      UCB는 \textbf{밴딧의 팔}(상태 없는 단일 선택) 하나하나에, UCT는
      \textbf{게임 트리의 매 노드마다} 독립적으로 적용되어 ``각 국면에서
      다음 수를 고르는'' 문제를 반복해서 푼다.}
    \item \textbf{2.} Ch6.1의 \kb{TD 오차}와 Ch10.3의
      \kb{Actor-Critic}은 어떤 관계인가?
      \vskip0.2em
      {\footnotesize $\to$ Actor-Critic의 \textbf{critic이 학습하는
      신호가 바로 TD 오차} — REINFORCE(10.2)가 에피소드 전체 리턴으로
      그래디언트 분산이 컸다면, AC는 그 리턴을 \textbf{TD 오차}(한
      스텝만 보고 계산한 어드밴티지 추정치)로 대체해 \textbf{분산을
      줄인다}. Ch11.3의 \textbf{GAE}는 이 TD 오차를 여러 스텝에
      걸쳐 지수가중평균한 \textbf{일반화}.}
  \end{itemize}
\end{frame}
